% Section 1.7, from lectures/L01/README.md: the goals, the self-check questions and the look
% ahead.

\section{Summary}
\label{c1:sec:review}

\needspace{6\baselineskip}
\subsection*{What you should be able to do now}
\begin{itemize}
  \item Derive a sum of products from a truth table, and then build and simulate the network in
    CircuitVerse.
  \item Apply De Morgan to move an inversion through a gate, and use that to realise a function in
    a gate set that lacks the operator you started from.
  \item Read and write a combinational VHDL module, an \code{entity} with \code{std_logic} ports
    and an \code{architecture} that drives its output, and say why VHDL keeps the two apart.
  \item Explain why a concurrent assignment describes a wire rather than a statement that runs
    once, and why a signal left sitting at \code{'U'} or \code{'X'} is a fault and not a value.
  \item Run a testbench handed out with an exercise and decide whether your module passed it.
\end{itemize}

\needspace{6\baselineskip}
\subsection*{Questions to test yourself with}
\begin{enumerate}
  \item What is the difference between an \code{entity} and an \code{architecture}, and why does
    VHDL separate them?
  \item Why is \code{x <= a or b;} not ``a statement that runs once''? What does it describe
    instead?
  \item In the brake assistant the driver's pedal is ORed in last, instead of going through the
    fault-handling logic. What breaks if you wire it the other way round?
\end{enumerate}

\needspace{6\baselineskip}
\subsection*{Looking ahead}
\Chapref{c2:ch} goes on with:
\begin{itemize}
  \item Karnaugh maps: finding the \emph{smallest} gate network, not just a correct one.
  \item \code{std_logic_vector}, \code{process} and \code{case}.
  \item Submodules, demonstrated on a two-digit hexadecimal display driven by eight switches.
  \item Self-checking testbenches: verifying your own VHDL without a board.
\end{itemize}
